\documentclass[10pt]{article}

\usepackage{microtype,amsmath,amssymb,amsthm,booktabs}
\usepackage[margin=0.86in]{geometry}
\usepackage[colorlinks=true,linkcolor=blue,urlcolor=blue,citecolor=blue]{hyperref}
\usepackage{bookmark}
\newtheorem{theorem}{Theorem}
\newtheorem{lemma}[theorem]{Lemma}
\newtheorem{corollary}[theorem]{Corollary}
\newcommand{\dd}{\mathrm d}
\newcommand{\R}{\mathbb R}
\setlength{\emergencystretch}{2em}
\title{Concentration Eigenvalues from Logarithmic Connection Coefficients}
\author{Kirtus Justus (Theoretical Physicist)\\
\small Port Coquitlam, British Columbia, Canada}
\date{October 2026}
\hypersetup{pdftitle={Concentration Eigenvalues from Logarithmic Connection Coefficients},pdfauthor={Kirtus Justus (Theoretical Physicist)}}
\usepackage{fontspec,unicode-math,mathtools}
\setmainfont{DejaVu Serif}\setsansfont{DejaVu Sans}\setmonofont{DejaVu Sans Mono}\setmathfont{Latin Modern Math}
\begin{document}
\maketitle
\begin{abstract}
For a finite sine kernel at arbitrary filling, we express each concentration
odds ratio exactly as a ratio of derivatives of logarithmic connection
coefficients for complementary Sturm--Liouville equations. A Lagrange
identity identifies those derivatives with positive norming integrals.
The formula includes the parity factor $(-1)^{N-1}$; the familiar minus sign
at half filling therefore requires even matrix size. Dividing out the finite
Jacobi characteristic polynomial gives an equivalent formula without
derivatives. We also give quantitative criteria converting relative
connection or solution-norm errors into concentration-logit errors.
These are finite identities and sufficient error criteria. Identification
with a formal quantum period, including a uniform spectral remainder,
is supplied at half filling on fixed interior intervals by the separate companion \emph{Uniform Bulk Modular Asymptotics and the Leading Wilson Entropy Law}. A general-filling uniform remainder remains a separate problem.
\end{abstract}

\section{Finite model and conventions}
Let $N\ge1$, $0<\theta<\pi$, $m_j=j-(N-1)/2$ for $0\le j<N$,
and $\alpha=(N^2-1)/4$. Define the real symmetric matrices
\begin{align}
(B_\theta)_{jk}&=
\begin{cases}
\theta/\pi,&j=k,\\
\sin[\theta(j-k)]/[\pi(j-k)],&j\ne k,
\end{cases} \label{eq:B}\\
(H_\theta)_{jj}&=\cos\theta\,m_j^2,\qquad
(H_\theta)_{j-1,j}=(H_\theta)_{j,j-1}=\frac{j(N-j)}2.
\label{eq:H}
\end{align}
Here $\theta=2\pi W$ in the bandwidth convention $0<W<1/2$.
The concentration formulation and its commuting tridiagonal operator are
classical; see Slepian \cite{slepian} and the explicit free-fermion formulation
in Eisler--Peschel \cite{eisler}. Our diagonal convention differs from a
frequently used one by a scalar shift. All normalizations needed below
are proved directly.

For $U_v(x)=\sum_j v_j e^{im_jx}$, Parseval's identity gives
\begin{equation}
v^*B_\theta v=\frac1{2\pi}\int_{-\theta}^{\theta}|U_v(x)|^2\,\dd x,
\qquad
\|v\|^2=\frac1{2\pi}\int_{-\theta}^{2\pi-\theta}|U_v(x)|^2\,\dd x.
\label{eq:parseval}
\end{equation}
Consequently $0<B_\theta<I$: a nonzero finite Fourier polynomial cannot
vanish on an open interval. A concentration eigenvalue is denoted $\zeta$,
and its logit is
\begin{equation}
\mathcal P=\log\frac{\zeta}{1-\zeta}=-\varepsilon.
\label{eq:logit}
\end{equation}
All logarithms are natural. No physical time or energy scale is assigned
to the differential-equation parameters by these definitions.

\begin{lemma}[Exact Fourier equation and commutation]
Put $p_\theta(x)=\cos x-\cos\theta$ and
$\mathcal L_\theta U=(p_\theta U')'+\alpha\cos x\,U$.
Then $\mathcal L_\theta U_v=U_{H_\theta v}$ and
$[H_\theta,B_\theta]=0$.
For each eigenpair $H_\theta v=\beta v$, the Fourier function satisfies
\begin{equation}
(p_\theta U')'+(\alpha\cos x-\beta)U=0,\qquad
U(x+2\pi)=\sigma U(x),\quad \sigma=(-1)^{N-1}.
\label{eq:ode}
\end{equation}
The eigenvalues of $H_\theta$ are simple, so its eigenvectors are also
eigenvectors of $B_\theta$.
\end{lemma}
\begin{proof}
For a Fourier monomial $e^{imx}$, the coefficients at frequencies $m+1,m-1$
are respectively $[\alpha-m(m+1)]/2$ and $[\alpha-m(m-1)]/2$.
The coefficient at $m$ is $\cos\theta\,m^2$. These are precisely
\eqref{eq:H}; the two outer coefficients vanish at the ends of the finite
frequency set. Integration by parts on $[-\theta,\theta]$ gives
$H_\theta B_\theta=B_\theta H_\theta$, since $p_\theta$ vanishes at
both endpoints and the Fourier functions and their derivatives are bounded.
The positive off-diagonal entries make $H_\theta$ an irreducible Jacobi
matrix, with simple spectrum; $N=1$ is immediate.
Finally $e^{2\pi i m_j}=\sigma$.
\end{proof}

\section{A normalized connection coefficient and its derivative}
For a complex spectral parameter $b$, let $u_\theta(x;b)$ be the solution of
\begin{equation}
(p_\theta u')'+(\alpha\cos x-b)u=0
\label{eq:param}
\end{equation}
analytic at $x=-\theta$, with $u_\theta(-\theta;b)=1$.
At either endpoint the indicial root is zero with multiplicity two.
The analytic solution with specified endpoint value is unique:
its Taylor recurrence divides the coefficient of degree $k$ by
$k^2\sin\theta\ne0$. In particular, an analytic solution that vanishes
at an endpoint vanishes identically.

Set $v_\theta(x;b)=u_\theta(-x;b)$, the analytic solution normalized to one
at the right endpoint. Define $C_\theta(b)$ by the local form
\begin{equation}
u_\theta(x;b)=C_\theta(b)v_\theta(x;b)\log(\theta-x)
+a_\theta(x;b),
\label{eq:C}
\end{equation}
where $a_\theta$ is analytic near $\theta$. This fixes the logarithmic
normalization. All derivatives of $C_\theta$ below are with respect to
its own argument $b$, holding $\alpha$ and $\theta$ fixed.

\begin{lemma}[Wronskian and norming identities]
The function $C_\theta$ is entire in $b$, and
\begin{equation}
C_\theta(b)=-\frac{2(1-\cos\theta)}{\sin\theta}
u_\theta(0;b)\,\partial_xu_\theta(0;b).
\label{eq:W}
\end{equation}
At a real zero $b_0$ of $C_\theta$, write $h=u_\theta(\theta;b_0)$.
Then $h\in\{-1,1\}$ and
\begin{equation}
M_\theta(b_0):=\int_{-\theta}^{\theta}u_\theta(x;b_0)^2\,\dd x
=-\sin\theta\,h\,C_\theta'(b_0)>0.
\label{eq:norm}
\end{equation}
In particular, every real zero of $C_\theta$ is simple.
\end{lemma}
\begin{proof}
The Wronskian $p_\theta(uv'-u'v)$ is constant. Its right-endpoint limit
is $\sin\theta\,C_\theta$, while its value at zero is
$-2(1-\cos\theta)u_\theta(0)\partial_xu_\theta(0)$.
This proves \eqref{eq:W}. The endpoint series is analytic in $b$;
continuation across compact subintervals preserves that dependence.
Thus both factors in \eqref{eq:W}, and hence $C_\theta$, are entire.

At a real zero the solution is real and analytic at both endpoints,
and $h\ne0$ by endpoint uniqueness. Reflection and uniqueness give
$u_\theta(-x;b_0)=h\,u_\theta(x;b_0)$; evaluation at the other endpoint
gives $h^2=1$.
Let $w=\partial_bu_\theta|_{b_0}$.
Differentiating \eqref{eq:param} and subtracting the two equations yields
\[
\bigl[p_\theta(uw'-u'w)\bigr]'=u^2.
\]
The left endpoint contributes zero. At the right endpoint,
$w=C_\theta'(b_0)v_\theta\log(\theta-x)+O(1)$, and the boundary
term is $-\sin\theta\,h\,C_\theta'(b_0)$.
Integration proves \eqref{eq:norm}, including its sign and nonvanishing.
\end{proof}

\section{Exact concentration formula at every filling}
\begin{theorem}[Complementary-band connection formula]
For every eigenmode $H_\theta v=\beta v$, $B_\theta v=\zeta v$, set
$\bar\theta=\pi-\theta$. Then
\begin{equation}
\boxed{\quad
\frac{\zeta}{1-\zeta}
=(-1)^{N-1}\frac{C_\theta'(\beta)}{C_{\bar\theta}'(-\beta)}.
\quad} \label{eq:main}
\end{equation}
Both denominator and numerator are nonzero. Equivalently,
\begin{equation}
\mathcal P=\log|C_\theta'(\beta)|-\log|C_{\bar\theta}'(-\beta)|.
\label{eq:logC}
\end{equation}
\end{theorem}
\begin{proof}
The entire Fourier solution of \eqref{eq:ode} is nonzero at both
$-\theta$ and $\theta$; otherwise endpoint uniqueness would make it
identically zero. Therefore
$u_\theta(x;\beta)=U(x)/U(-\theta)$ is the real normalized regular
solution, $C_\theta(\beta)=0$, and $U(\theta)=hU(-\theta)$ with $h=\pm1$.

Translate the complementary arc by $x=\pi+y$. Since
$p_\theta(\pi+y)=-p_{\bar\theta}(y)$, the translated equation is
\eqref{eq:param} with $(\theta,b)$ replaced by $(\bar\theta,-\beta)$.
Its near-normalized solution is $U(\pi+y)/U(\theta)$, and its far value is
\[
\bar h=\frac{U(2\pi-\theta)}{U(\theta)}
=\frac{\sigma}{h}=\sigma h.
\]
The two arc normalizations have equal absolute value.
Equation \eqref{eq:parseval} therefore gives
$\zeta/(1-\zeta)=M_\theta(\beta)/M_{\bar\theta}(-\beta)$.
Apply \eqref{eq:norm}, using $\sin\bar\theta=\sin\theta$ and
$h/\bar h=\sigma$, to obtain \eqref{eq:main}.
\end{proof}

At half filling $\theta=\pi/2$, translating $s=x+\pi/2$ gives
\[
(\sin s\,u')'+(\alpha\sin s-b)u=0,\qquad 0<s<\pi.
\]
If $J$ has bonds $j(N-j)/N$ and zero diagonal, then
$H_{\pi/2}=NJ/2$ and $\beta=N\lambda/2$.
For even $N$, \eqref{eq:main} is precisely
$\zeta/(1-\zeta)=-C'(\beta)/C'(-\beta)$.
For odd $N$, the sign is positive. For example, at the central mode
$\beta=0$ one has $\zeta=1/2$; a fixed minus sign would give $-1$
instead of $1$. More explicitly, for $N=2m+1$ the regular central
solution is $P_m(\cos s)$ and
\begin{equation}
C'(0)=-(-1)^m\int_0^\pi P_m(\cos s)^2\,\dd s\ne0.
\label{eq:central}
\end{equation}
The norming integral here uses $\dd s$, not $\sin s\,\dd s$.

\begin{corollary}[Removal of the finite characteristic polynomial]
Let $D_\theta(b)=\det(bI-H_\theta)$ and
$R_\theta(b)=C_\theta(b)/D_\theta(b)$, with removable values filled in.
Then $R_\theta$ is entire, is nonzero at every eigenvalue of $H_\theta$,
and at each such eigenvalue
\begin{equation}
\frac{\zeta}{1-\zeta}
=\frac{R_\theta(\beta)}{R_{\bar\theta}(-\beta)}
=(-1)^N\lim_{b\to\beta}\frac{C_\theta(b)}{C_{\bar\theta}(-b)}.
\label{eq:deflate}
\end{equation}
\end{corollary}
\begin{proof}
The $N$ distinct zeros of $D_\theta$ are simple zeros of $C_\theta$,
so division removes only simple zeros.
For $E_{jj}=(-1)^j$, $H_{\bar\theta}=-EH_\theta E$, whence
$D_{\bar\theta}(-b)=(-1)^ND_\theta(b)$ and
$D_{\bar\theta}'(-\beta)=\sigma D_\theta'(\beta)$.
Substitute $C_\theta'(\beta)=D_\theta'(\beta)R_\theta(\beta)$
into \eqref{eq:main}. The last expression follows by taking the limit.
\end{proof}
This does not identify the full zero set of $C_\theta$ with the finite
matrix spectrum. Nor is \eqref{eq:deflate} asserted to be a concentration
formula away from those spectral points.

\section{Quantitative routes to a period remainder}
The exact formula isolates the missing asymptotic task without solving it.
Two elementary stability statements specify sufficient accuracy.

\begin{theorem}[Relative-error criteria]
At an exact eigenvalue $\beta$, put
$d_+=C_\theta'(\beta)$, $d_-=C_{\bar\theta}'(-\beta)$.
If nonzero approximations $\widehat d_\pm$ satisfy
$|\widehat d_\pm/d_\pm-1|\le q_\pm<1$, then
\begin{equation}
\left|\mathcal P-\log\left|\frac{\widehat d_+}{\widehat d_-}\right|\right|
\le-\log(1-q_+)-\log(1-q_-).
\label{eq:relative}
\end{equation}
Alternatively let $u_\pm$ be the two near-normalized regular arc solutions
in the theorem above. For approximants $f_\pm$ with $F_\pm=\|f_\pm\|_2>0$,
suppose
$\|u_\pm-f_\pm\|_2\le\delta_\pm F_\pm$, $0\le\delta_\pm<1$.
Writing $\widehat{\mathcal P}=\log(F_+^2/F_-^2)$ gives
\begin{align}
2\log(1-\delta_+)-2\log(1+\delta_-)
&\le\mathcal P-\widehat{\mathcal P}\nonumber\\
&\le2\log(1+\delta_+)-2\log(1-\delta_-).
\label{eq:L2}
\end{align}
\end{theorem}
\begin{proof}
The reverse triangle inequality bounds
$|\widehat d_\pm/d_\pm|$ between $1-q_\pm$ and $1+q_\pm$.
Apply \eqref{eq:logC} and
$\log(1+q)\le-\log(1-q)$.
For the second assertion,
$(1-\delta_\pm)F_\pm\le\|u_\pm\|_2\le(1+\delta_\pm)F_\pm$.
Square, take the ratio, and use the norm formula in the proof of
\eqref{eq:main}.
\end{proof}

These bounds do not contain a factor $1/[\zeta(1-\zeta)]$:
relative norming data control the logit directly even near concentration
endpoints. But the hypotheses require the correct endpoint normalization
and the actual spectral parameter. A small residual for a homogeneous
ODE alone cannot set the normalization: multiplying a solution by a
constant preserves that equation and changes its squared norm.
Use of an approximate $\beta$ also needs an independent eigenvalue and
parameter-sensitivity bound.

For the half-filled model with even $N$, set $t=2\beta/N^2$ and define the finite
formal-period truncation
\begin{equation}
P_N^{[5]}(t)=N\eta(t)-\frac{\varphi_1(t)}N
-\frac{\varphi_3(t)}{N^3}-\frac{\varphi_5(t)}{N^5}.
\label{eq:period}
\end{equation}
Here $\eta(t)=\pi t\,{}_3F_2(1/4,1/2,3/4;1,3/2;4t^2)$ and the
$\varphi_k$ are the explicitly defined formal cycle coefficients
in \cite{periods}. For a compact $K\subset(-1/2,0)\cup(0,1/2)$, sufficient
conditions for the desired spectral theorem are uniform
$q_\pm=O_K(N^{-7})$ in \eqref{eq:relative} and
\[
\log|\widehat d_+/\widehat d_-|
=P_N^{[5]}(t)+O_K(N^{-7})
\]
at the modes with $t\in K$. Equation \eqref{eq:L2} provides an alternative
through normalized solution norms. Either set of conditions implies
\begin{equation}
\left|e-\frac{\varphi_1(t)}N-\frac{\varphi_3(t)}{N^3}
-\frac{\varphi_5(t)}{N^5}\right|\le C_KN^{-7},
\qquad e=-\mathcal P+N\eta(t).
\label{eq:target}
\end{equation}
The implication is proved; the stated uniform approximation hypotheses
have not been established here. A formal all-orders series also needs a
defined summation before it can be called an exact period function.

\section{Verification and scope}
The identities above have analytic proofs independent of the numerical
checks. A separate computation obtains $\beta$ and $\zeta$ from the
finite matrices and computes $C$ and $C'$ from a parameter-differentiated
Frobenius recurrence and the midpoint Wronskian \eqref{eq:W}.
The concentration eigenvalue is not used in that connection calculation.
At 90 decimal digits and Taylor degree $380$, 74 cases at
$\theta/\pi=1/4,1/2,3/4$, with selected sizes $1\le N\le25$, gave a
largest relative odds discrepancy below $3.43\times10^{-73}$.
These comprised 48 half-filled cases and 26 other-filling cases.
Four repetitions at 120 digits and degree $500$ gave discrepancies below
$9.77\times10^{-114}$; separate complex-step checks tested the parameter
derivative. Exact low-order solutions also fix the logarithmic normalization.
Ordinary multiprecision agreement is not interval validation.

These results refine the finite reference problem in
\emph{From Prephysical Hypotheses to Physical Models}.
They do not establish the uniform period expansion, the full Wilson
relative-entropy law, or a Yang--Mills mass-gap theorem.
The exact complementary-band formula and the sufficient error criteria
are the established conclusions here. The whole-line construction in \emph{Uniform Bulk Modular Asymptotics
and the Leading Wilson Entropy Law} proves the half-filled
fixed-interior expansion. For interval restrictions of the specified
infinite-chain Wilson--Dirac projectors with \(m=\mu/L\) and fixed
\(\mu>0\), Theorem 1 of \emph{Finite Part of Wilson--Dirac Relative
Entropy} determines the finite part with remainder \(O_\mu(1/\log L)\).
Theorem 3 of \emph{Uniform Modular Control and Nonlinear Convergence for
Wilson--Dirac Intervals} also proves \(\|E_N\|\le7\) for even \(N\ge2\).
These conclusions do not establish the sufficient endpoint approximation
hypotheses stated in this article, a general-filling expansion, or a
moving-edge expansion.

\paragraph{Acknowledgments and provenance.}
The even-$N$, half-filled connection formula was supplied in research
correspondence prepared with Claude (Anthropic).
The present derivation, parity clarification, complementary-band extension,
and checks were developed with ChatGPT (OpenAI).
The wider research program has also used Gemini (Google),
Grok (SpaceXAI), and DeepSeek (High-Flyer); no specific derivation in this
paper is attributed to those systems. Assistance does not imply
independent external peer review or endorsement. No priority claim is made.

\begin{thebibliography}{9}
\bibitem{slepian} D. Slepian, Prolate spheroidal wave functions, Fourier
analysis, and uncertainty---V: The discrete case,
\emph{Bell Syst. Tech. J.} \textbf{57} (1978), 1371--1430.
Original source cited following \cite{eisler}; its full text was not
independently re-read for this note.
\bibitem{eisler} V. Eisler and I. Peschel,
Analytical results for the entanglement Hamiltonian of a free-fermion chain,
\emph{J. Phys. A} \textbf{50} (2017), 284003,
\href{https://arxiv.org/abs/1703.08126}{arXiv:1703.08126},
especially equations (3), (6)--(10).
\bibitem{periods} K. Justus,
\emph{Higher-Order Modular Periods for the Half-Filled Sine Kernel},
research manuscript, October 2026.
\end{thebibliography}
\end{document}
